% paper1_indicators.tex -- Paper 1. Typeset alone with this file, or inside
% the volume through main.tex.
\documentclass[../main.tex]{subfiles}
\begin{document}
\begin{paper}{Selecting a Minimum Dataset of Soil Health Indicators for Mixed Cropping Systems}{A minimum dataset of soil health indicators}
\label{paper:indicators}

\begin{paperabstract}
Soil health cannot be measured directly; it has to be inferred from a small
number of soil properties that respond to management and that relate to the
functions a farmer depends on. This paper describes how a minimum dataset of
six indicators was chosen for a multi-site trial network covering three
climatic zones, how each indicator is converted to a common score between 0
and 1, and how the scores are combined into a single soil health index. The
dataset spans chemical, physical and biological properties so that no single
group of processes dominates the index. The approach is applied to synthetic
trial data and is intended as a worked example of indicator selection and
scoring, not as a recommendation for any real location.
\end{paperabstract}
\paperkeywords{soil health; soil quality; minimum dataset; indicator scoring; soil health index}

\section{Introduction}

Interest in the condition of agricultural soils has grown steadily since
soil quality was defined as the capacity of a soil to function within
ecosystem boundaries, sustaining biological productivity, environmental
quality and plant and animal health \parencite{karlen1997}. The term
\emph{soil health} places more weight on the soil as a living system, and on
the organisms whose activity drives nutrient cycling and the formation of
soil structure \parencite{doran2000}. The two terms are now often used
interchangeably, although the concept of soil health has broadened to include
links with human health and with the services that soils provide beyond the
farm \parencite{lehmann2020}.

Because soil health is a property of the whole system, it cannot be read from
a single instrument. It is inferred instead from indicators: measurable soil
properties that are sensitive to management, related to a soil function, and
practical to measure. The number of candidate indicators is large, and a
review of assessment schemes found that the choice between them, and the way
measured values are interpreted, is still inconsistent from one programme to
the next \parencite{bunemann2018}. A common response is to select a
\emph{minimum dataset}, the smallest set of indicators that together describe
the functions of interest without redundancy.

Selecting indicators is only the first of three steps. Each measured value
must then be interpreted, usually by transforming it to a unitless score that
expresses how well the soil performs relative to a target, and the scores must
finally be combined into an index. The Soil Management Assessment Framework
formalised these three steps and showed that scoring curves which rise, fall
or peak with the measured value can represent the different ways in which an
indicator relates to soil function \parencite{andrews2004}.

This paper applies the same three steps to the trial network used throughout
this volume. Its aims are to state the criteria by which six indicators were
chosen, to define the scoring function for each, and to show how the resulting
index is calculated, so that the later papers can report soil health on a
common basis. Soil organic carbon is determined by wet oxidation
\parencite{walkley1934} and extractable phosphorus by the Mehlich~3 method
\parencite{mehlich1984}; the remaining indicators and all laboratory procedures
are described under Materials and Methods.

\input{\subfix{methods/paper1_methods}}

\end{paper}
\end{document}
